\documentclass{article}
\title{ECE 341 -- Homework 3}
\usepackage{graphicx}
\usepackage[margin=1in]{geometry}
\usepackage{float}
\usepackage{pdfpages}
\usepackage{fancyhdr}
\usepackage{enumitem}
\usepackage{pdfpages}
\usepackage{amsmath}
\usepackage{amssymb}
\pagestyle{fancy}
\fancyhf{}
\fancyhead[L]{Gianni Avilan-Losee}
\fancyhead[C]{ECE 341}
\fancyhead[R]{Homework 4}
\fancyfoot[C]{\thepage}
\usepackage{microtype}
\usepackage{textcomp, gensymb}
\author{Gianni Avilan-Losee}
\begin{document}
\includepdf{homework-4_cover-page.pdf}
\begin{enumerate}[label=\thesubsubsection]
\setcounter{section}{4}
\setcounter{subsection}{1}
\setcounter{subsubsection}{1}
\item 
	\begin{enumerate}[label=\alph*)]
		\item First, we may split the given function into two integrals, on the interval where the given unit functions resolve to non-zero values, as
			\[
				\int_{0}^{\infty}u(t)e^{-st} \ dt-\int_{1}^{\infty}u(t-1)e^{-st} \ dt = \int_{0}^{\infty}e^{-st} \ dt - \int_{1}^{\infty}e^{-st} \ dt.
			\]
			Then, evaluating the integrals,
			\begin{align*}
				=& \ \frac{-e^{-st}}{s} \ \Bigg\rvert_{0}^{\infty} + \frac{e^{-st}}{s} \ \Bigg\rvert_{1}^{\infty} \\ \\
				=& \left[\frac{-e^{-s \infty}}{s} + \frac{1}{s}\right] + \left[\frac{e^{-s\infty}}{s}-\frac{e^{-s}}{s}\right] \\ \\
				=& \ \frac{1-e^{-s}}{s}.
			\end{align*}
			The given function converges only when \(\sigma > 0\), as the imaginary part of \(s\) will continue to oscillate indefinitely (an everlasting sinusoid). Therefore, we may assume that the ROC for this function is \(\Re{(s)} > 0\).
			\addtocounter{enumii}{2}
		\item Again, splitting the given function into two integrals,
			\[
				\int_{0}^{\infty}e^{-2t}e^{-st}u(t) \ dt - 2 \int_{0}^{\infty}e^{-t}e^{-st}u(t) \ dt = \int_{0}^{\infty}-e^{-t(s+2)} \ dt + 2\int_{0}^{\infty}e^{-t(s+1)} \ dt. 
			\]
			Then, evaluating the integrals,
			\begin{align*}
				&= \ -\frac{e^{-2t}e^{-st}}{s+2} \ \Bigg\rvert_{0}^{\infty} + \frac{2e^{-t}e^{-st}}{s+1} \ \Bigg\rvert_{0}^{\infty} \\ \\
				&=\frac{1}{s+2} + \frac{2}{s+1}.
			\end{align*}
			The region of convergence for this function depends on both \(\sigma\) and \(j\omega\), the real and imaginary parts of \(s\). Under the condition where \(\sigma = -2\) and \(\omega = 0\), the function will \textit{not} converge, as well as under the condition where \(\sigma = -1\) and \(\omega = 0\). Therefore, the function will converge as long as \(s \neq -2\) and \(s \neq -1\) (implying an imaginary part of \(0\) in both cases).
	\end{enumerate}
	\setcounter{subsubsection}{4}
	\item
		\begin{enumerate}[label=\alph*)]
			\item Using partial fraction decomposition, 
				\begin{align*}
				 &\frac{2s+5}{s^2+5s+6} = \frac{2s+5}{(s+2)(s+3)} = \frac{A}{(s+2)}+\frac{B}{(s+3)} \\ \\
				 &2s+5 = A(s+3) + B(s+2) = sA + 3A +sB + 2B \\ \\
				 &2s+5 = s(A+B) +3A + 2B.
				\end{align*}
				Solving for \(A\) and \(B\), \(A+B = 2\) and \(3A + 2B = 5\), so
				\[
					\begin{bmatrix}
						1 & 1 & 2 \\ 3 & 2 & 5
					\end{bmatrix} \to 
					\begin{bmatrix}
						1 & 0 & 1 \\ 0 & 1 & 1
					\end{bmatrix},
				\]
				and so \(A = 1\) and \(B=1\).
				\\[10pt]
				Therefore, 
		\[
		  \frac{2s+5}{s^2+5s+6} = \frac{1}{(s+2)} + \frac{1}{(s+3)},
		\]
		and the inverse Laplace transform of the function may be found using Table 4.1, specifically the unilateral transform pair \(1/(s-\lambda) \Leftrightarrow e^{\lambda t}u(t)\),
		\begin{align*}
			&\mathcal{L}^{-1}\left\{\frac{1}{s+2} + \frac{2}{s+1}\right\} = \mathcal{L}^{-1}\left\{\frac{1}{s+2}\right\} + \mathcal{L}^{-1}\left\{\frac{1}{s+3}\right\} \\ \\
			=& \ [e^{-2t} + e^{-3t}]u(t).
		\end{align*}
  \end{enumerate}
	\setcounter{subsection}{2}
	\setcounter{subsubsection}{2}
	\item
		\begin{enumerate}[label=\alph*)]
			\item Using the Laplace transform pair \(u(t) \Leftrightarrow 1/s\) and the time-shifting property of the Laplace transform, \(x(t-t_0) \Leftrightarrow X(s)e^{-st_0}\),
				\begin{align*}
					&\mathcal{L}\left\{u(t) - u(t-1)\right\} = \mathcal{L}\left\{u(t)\right\} - \mathcal{L}\left\{u(t-1)\right\} \\
					=& \frac{1}{s} - \frac{e^{-s}}{s} = \frac{1-e^{-s}}{s}.\\
				\end{align*}
			\item Likewise, using the time-shifting property of the Laplace transform along with the Laplace transform pair \(e^{\lambda t}u(t) \Leftrightarrow 1/(s-\lambda)\),
				\[
					\mathcal{L}\left\{e^{-(t-\tau)}u(t-\tau)\right\} = \frac{e^{-s\tau}}{s+t}.
				\]
		\end{enumerate}
		\setcounter{subsection}{3}
		\setcounter{subsubsection}{1}
		\item
			\begin{enumerate}[label=\alph*)]
			  \addtocounter{enumii}{1}	
				\item Converting to differential form, 
					\[
						\ddot{y}(t) + 4\dot{y}(t) + 4y(t) = \frac{d}{dt}e^{-t}u(t) + e^{-t}u(t)
					\]
					with initial conditions
					\[
						y(0^-) = \dot{y}(0^-) = 0,
					\]
					the time differentiation property of the Laplace transform may be used along with the transform pair \(e^{\lambda t}u(t) \Leftrightarrow 1/(s-\lambda)\) to find
					\begin{align*}
					&\mathcal{L}\left\{\ddot{y}(t) + 4\dot{y}(t) + 4y(t)\right\} = \mathcal{L}\left\{\frac{d}{dt}e^{-t}u(t) + e^{-t}u(t)\right\} \\
						=&\left[s^2Y(s) - sy(0^-) - \dot{y}(0^-)\right] + 4\left[sY(s) - y(0^-)\right] + 4Y(s) = sX(s) - x(0^-) + \frac{1}{s+1}\\
						=& \ s^2Y(s) + 4sY(s) + 4Y(s) = \frac{s+1}{s+1} = 1\\
						=& \ Y(s)\left[s^2+4s+4\right] = 1.
					\end{align*}
					Solving algebraically for \(Y(s)\),
					\[
					  Y(s) = \frac{1}{s^2 + 4s + 4} = \frac{1}{(s+2)^2},
					\]
					which may be solved directly using the Laplace transform pair \(te^{\lambda t}u(t) \Leftrightarrow 1/(s-\lambda)^2\), as
					\[
						\mathcal{L}^-1\left\{\frac{1}{(s+2)^2}\right\} = te^{-2t}u(t).
					\]
			\end{enumerate}
			\setcounter{subsection}{4}
			\setcounter{subsubsection}{7}
		\item At \(t<0\), before the switch has opened, we may derive the initial conditions of the circuit as 
			\begin{align*}
				&v_s(0^-) = 0^{\mathrm{V}} \\
				&y(0^-) = 10^{\mathrm{A}} \\
				&y_2(0^-) = 0^{\mathrm{A}}.
			\end{align*}
			After the switch opens, at \(t > 0\), we may derive the voltage across each component using the voltage equations for inductors, resistors, and capacitors, 
			\begin{align*}
				&v_R(t) = Ri(t) \\
				&v_L(t) = L \frac{di_L(t)}{dt} \\
				&v_C(t) = \frac{1}{C}\int i(t) + v_0,
			\end{align*}
			along with the unilateral Laplace transform pairs 
			\begin{align*}
				&a \Leftrightarrow \frac{a}{s} \\
				&x(t) = X(s) \\
				&\frac{dx(t)}{dt} \Leftrightarrow sX(s) + x(0^-) \\
				&\int_{0^-}^{t}x(\tau)d\tau \Leftrightarrow \frac{1}{s}X(s) \\
			\end{align*}
			to derive the loop equation (KVL) for the circuit at \(t > 0\),
			\begin{align*}
				\frac{10}{s} &= 2sY(s) - 10 + sY(s) + \frac{1}{s}Y(s) + Y(s) \\
										 &=Y(s) \left[3s + \frac{1}{s} + 1\right] - 10, 
			\end{align*}
			so that 
			\[
				Y(s) = \frac{10s + 10}{3s^2+s+1}.
			\]
			Transforming \(Y(s)\) back to the time domain \textit{could} be done using partial fraction expansion, but a simpler method is to complete the square in the denominator and use the transform pairs
			\begin{align*}
				&e^{-at}\cos{bt}u(t) \Leftrightarrow \frac{s+a}{(s+a)^2+b^2} \\
				&e^{-at}\sin{bt}u(t) \Leftrightarrow \frac{b}{(s^2+b^2)} \\
			\end{align*}.
			Mathematically,
			\begin{align*}
				Y(s) &= \frac{10}{3}\left[\frac{(s+1)}{(s^2+s/3+1/3)}\right]\\
						 &= \frac{10}{3}\left[\frac{s+1/6+5/6}{(s+1/6)^2 + 11/36}\right] \\
						 &= \frac{10}{3}\left[\frac{s+1/6}{(s+1/6)^2} + \frac{5/6}{(s+1/6)^2+11/36}\right].
			\end{align*}
			Then, using the tranform pairs above, 
			\[
				y(t) = \frac{10}{3}e^{\frac{-t}{6}} \left[\cos{(\frac{\sqrt{11}}{6}t)} + \frac{5}{\sqrt{11}}\sin{(\frac{\sqrt{11}}{6}t)}\right]u(t) = y_1(t).
			\]
			To find the voltage across the series inductor and capacitor, we may use their s-domain expressions, sum them, and convert back to the time domain. For their voltages, 
			\[
			  V(s) = sY(s) + \frac{1}{s}Y(s).
			\]
			Substituting \(Y(s)\),
			\[
			  V(s) = \frac{10s^2 + 10s}{3s^2+s+1} + \frac{10s+10}{3s^3+s+1}.
			\]
			Using similar methods as before to convert to the time domain,
			\[
				v(t) = \frac{10}{3}\delta(t) + 10u(t) + e^{\frac{-t}{6}}\left[\frac{-70}{9}\cos{(\frac{\sqrt{11}}{6}t)} + \frac{50}{9\sqrt{11}}\sin{(\frac{\sqrt{11}}{6}t)}\right]u(t).
			\]
			\setcounter{subsection}{8}
			\setcounter{subsubsection}{2}
		\item Note that for the given transfer function \(H(s)\), we may assume the system beign described is causal, since \(H(s)\) is \textit{proper}, and so the system response \(Y(s)\) may also be assumed to be causal, applying only under the condition \(t > 0\).
			\begin{enumerate}[label=\alph*)]
				\item Converting the input \(x(t) = 5\cos{(2t+30\degree)}\) to the s-domain using a unilateral Laplace transform, based on the transform pair
					\[
						re^{-at}\cos{(bt+\theta)}u(t) \Leftrightarrow \frac{(r\cos{\theta})s + (ar\cos{\theta}-br\sin{\theta})}{s^2 + 2as +(a^2+b^2)},
					\]
					we find
					\[
						X(s) = \frac{5s\cos{30\degree - 10\sin{30\degree}}}{s^2  + 2^2} = \frac{5s \sqrt{3}/2 - 5}{s^2 + 4}.
					\]
					Multiplying the transfer function by the input in the s-domain (equivalent to convolution in the time domain), and expressing both \(H(s)\) and \(X(s)\) in root and pole form for simplicity,
					\[
						Y(s) = \left[\frac{(s+2)}{(s+1)(s+4)}\right]\left[\frac{5(s\frac{\sqrt{3}}{2} - 1)}{(s+j2)(s-j2)}\right] = 5\left[\frac{0.866s^2+0.7321s - 2}{(s+1)(s+4)(s+j2)(s-j2)}\right]. 
					\]
				\item Once again, we will convert the input \(x(t) = 10\sin{(2t+45\degree)}\) to the s-domain by first using the identity
					\[
						\sin{x} = \cos{x-90\degree}
					\]
					to avoid using a different transform pair than in part a), so that 
					\[
						x(t) = 10\cos{(2t - 45\degree)}.
					\]
					Converting to the s-domain, 
					\[
						X(s) = \frac{10s\cos{-45\degree} - 20\sin{-45\degree}}{s^2 + 2^2} = \frac{10s/\sqrt{2} + 20/\sqrt{2}}{s^2+4}.
					\]
					Multiplying by the transfer function to find \(Y(s)\),
					\[
						Y(s) = \left[\frac{(s+2)}{(s+1)(s+4)}\right]\left[\frac{10(s+2)}{\sqrt{2}(s+j2)(s-j2)}\right] = 7.07\left[\frac{s^2+4s+16}{(s+1)(s+4)(s+j2)(s-j2)}\right].
					\]
				\item Finally, we will convert the input \(x(t) = 10\cos{(3t+40\degree)}\) to the s-domain using the same Laplace transform pair as before, as
					\[
						X(s) = \frac{10s\cos{40\degree} - 30\sin{40\degree}}{s^2 + 9} = \frac{7.66s-19.28}{(s^2+9)}.
					\]
					Then, multiplying by \(H(s)\) to find \(Y(s)\),
					\[
						Y(s) = \left[\frac{(s+2)}{(s+1)(s+4)}\right]\left[\frac{7.66s - 19.28}{(s+j2)(s-j2)}\right] = 7.66\left[\frac{s^2-0.52s+5.03}{(s+1)(s+4)(s+j2)(s-j2)}\right].
					\]
			\end{enumerate}
\end{enumerate}
\end{document}
